\begin{document}

\loomcover
  {信息论}
  {第三讲 · 信道容量\\[2pt]{\small\color{selvage}Oxford B8.4 · 2024--25}}
  {fanzuituanhuo}
  {\today}

\section*{一条主线与知识地图}
\phantomsection
\addcontentsline{toc}{section}{一条主线与知识地图}
\warmth{0}

\begin{strand}
信道无非是一个条件分布 $p(y\mid x)$。这个分布我们无法改动，但可以选择各个输入被
使用的频率。\keyword{容量}就是这种选择权能换来的最大互信息：
$C=\max_{p(x)}\MI(X;Y)$。本讲所有例子都是同一个两步动作的实例——先拆开
$\MI(X;Y)=\Ent(Y)-\Ent(Y\mid X)$，再在信道允许的范围内把输出的不确定性做到最大。
\end{strand}

\block{目前的知识地图}
\noindent\begin{tabularx}{\linewidth}{@{}c l L@{}}
\toprule
{\color{indigo}站} & {\color{indigo}对象} & \multicolumn{1}{l}{\color{madder}要回答的问题}\\
\midrule
1 & $(\X,p(y\mid x),\Y)$ & 描述一个信道需要哪些数据？\\
2 & $p(y^n\mid x^n)=\prod_ip(y_i\mid x_i)$ & “无记忆”排除了什么？\\
3 & $C=\max_{p(x)}\MI(X;Y)$ & 使用信道一次最多能传多少信息？\\
4 & $C=1$（无噪二元信道） & 一条完美的比特管道是什么样子？\\
5 & $\Ent(X\mid Y)=0$ & 什么时候噪声完全不付代价？\\
6 & $C=1-h(p)$（BSC） & 比特被翻转要付多少代价？\\
7 & $C=1-\alpha$（BEC） & 比特被丢失要付多少代价？\\
8 & $C=\log|\Y|-\Ent(r)$ & 对称性能白送我们什么？\\
9 & $0\le C\le\log\min(|\X|,|\Y|)$ & $C$ 究竟可能取哪些值？\\
\bottomrule
\end{tabularx}

\block{全文约定}
字母表 $\X$ 与 $\Y$ 都有限。所有对数取 $2$ 为底，因此熵与容量的单位是
比特每次信道使用。记二元熵函数 $h(p)=-p\log p-(1-p)\log(1-p)$，并用 $\Ent(r)$
表示概率向量 $r$ 的熵。输入分布 $p(x)$ 由我们选择，转移概率 $p(y\mid x)$ 不由我们
选择。

\clearpage
{\small\tableofcontents}

\clearpage
\section*{第三讲 · 信道容量}
\phantomsection
\addcontentsline{toc}{section}{第三讲 · 信道容量}
\begin{strand}
先固定“信道是什么”以及“无记忆”的含义；再把容量定义为对输入分布的最大化，并在四个
标准例子里算出来：无噪二元信道、输出互不重叠的噪声信道、二元对称信道、二元删除
信道。最后提炼出让这些计算变得容易的结构性理由——对称性，并记下 $C$ 的一般性质。
\end{strand}

% ===========================================================================
\section{信道即条件分布}\label{sec:channel-model}
\warmth{0}\quad\whisper{信道就是一张条件概率表；容量是我们能把它用到多好。}

\begin{strand}
离散信道的全部信息都装在一个随机矩阵里。矩阵由物理给定，剩下自由的只有喂入输入
时所用的分布。容量就是这份自由度的价值。
\end{strand}

\trigger{当一个符号被干扰的方式只依赖该符号本身时，就用这个模型。}

\block{信道的数据}

\begin{definition}[离散信道]\label{def:discrete-channel}
一个\keyword{离散信道}由输入字母表 $\X$、输出字母表 $\Y$ 以及概率转移矩阵
$p(y\mid x)$ 组成，其中 $p(y\mid x)$ 表示发送输入符号 $x$ 时观测到输出符号 $y$
的概率。它是随机矩阵，即
\[
  p(y\mid x)\ge0
  \quad\text{对一切 }x,y,
  \qquad
  \sum_{y\in\Y}p(y\mid x)=\answerin[1.2cm]{1}
  \quad\text{对每个 }x\in\X.
\]
\studentrecall{转移矩阵的每一\emph{行}都是 $\Y$ 上的概率分布，行和为 $1$。}
之所以称为\keyword{离散}，是因为两个字母表都有限。
\end{definition}

\begin{definition}[无记忆]\label{def:memoryless}
若信道的输出只与\answerin[3.2cm]{当前时刻的输入}有关，而与更早的输入或输出无关，
则称信道\keyword{无记忆}。等价地，对每个 $n$ 以及每个 $x^n\in\X^n$、
$y^n\in\Y^n$，
\[
  p(y^n\mid x^n)=\prod_{i=1}^{n} \answerin[2.4cm]{p(y_i\mid x_i)}.
\]
\studentrecall{只有当前输入起作用，于是分组分布分解为单次使用的 $p(y_i\mid x_i)$ 之积。}
既离散又无记忆的信道称为\keyword{离散无记忆信道}（DMC）。
\end{definition}

\begin{remark}
乘积公式悄悄用到了两个假设：信道不记得过去的输入（无记忆），以及输入 $x_i$ 不允许
依赖过去的输出 $y^{i-1}$（无反馈）。写成上式需要两者同时成立，本讲的设定下二者都
成立。
\end{remark}

\block{容量}

\begin{definition}[信息信道容量]\label{def:information-capacity}
离散无记忆信道的（“信息”）\keyword{信道容量}为
\[
  C=\max_{p(x)}\MI(X;Y),
\]
其中最大值取遍所有可能的\answerin[4.2cm]{输入分布 $p(x)$}。
\studentrecall{对输入分布 $p(x)$ 取最大；转移矩阵 $p(y\mid x)$ 是固定的。}
联合分布始终是 $p(x,y)=p(x)\,p(y\mid x)$，所以选定 $p(x)$ 就同时决定了输出分布
$p(y)=\sum_xp(x)p(y\mid x)$ 与 $\MI(X;Y)$ 的取值。
\end{definition}

\begin{remark}
容量还有一个\keyword{操作性}定义：能以趋于 $0$ 的错误概率传输信息的最大速率
（比特每次信道使用）。Shannon 第二定理（信道编码定理）说明信息信道容量与操作性
信道容量相等，故“information”一词常被省去。本讲计算 $\max_{p(x)}\MI(X;Y)$ 这一侧，
编码定理留给下一讲。
\end{remark}

\block{每个例子都在用的同一个动作}

由于 $\MI(X;Y)=\Ent(Y)-\Ent(Y\mid X)$ 且
\[
  \Ent(Y\mid X)=\sum_{x}p(x)\,\Ent(Y\mid X=x),
\]
第二项只是转移矩阵各行熵（这些行熵是\emph{固定}的）的加权平均。因此最大化
$\MI(X;Y)$ 就是在把平均行熵压小的同时把 $\Ent(Y)$ 做大。当所有行的熵\emph{相同}
时，第二项根本不依赖 $p(x)$，只需最大化 $\Ent(Y)$——这正是第
\ref{sec:bsc} 节与第 \ref{sec:symmetric} 节所用的捷径。

\begin{yourturn}
写出本讲反复使用的 $\MI(X;Y)$ 两项分解，并说明哪一项由信道固定、哪一项由输入分布
调动。
\studentrecall{$\MI=\Ent(Y)-\Ent(Y\mid X)$；$p(y\mid x)$ 的各行固定了条件熵，$p(x)$ 调动权重与 $\Ent(Y)$。}
\ifshowanswers\else\workspace[2]\fi
\answerprose{$\MI(X;Y)=\Ent(Y)-\Ent(Y\mid X)$，其中
$\Ent(Y\mid X)=\sum_xp(x)\Ent(Y\mid X=x)$。信道固定了每个行熵
$\Ent(Y\mid X=x)$；输入分布决定这些行熵的权重，并通过
$p(y)=\sum_xp(x)p(y\mid x)$ 控制 $\Ent(Y)$。}
\end{yourturn}

\recall{$p(y^n\mid x^n)=\prod_ip(y_i\mid x_i)$ 里隐藏了哪两个假设？}


\clearpage
% ===========================================================================
\section{两个噪声不要钱的信道}\label{sec:first-examples}
\warmth{0}\quad\whisper{只要噪声从不让两个输入看起来一样，它就是无害的。}

\begin{strand}
先看两个极端。信道无噪时，每次使用传一个无差错比特是显然的。信道有噪但不同输入
产生的输出集合互不相交时，接收端仍能精确恢复输入——容量还是等于一条干净管道。
真正消耗容量的不是随机性，而是\emph{可混淆性}。
\end{strand}

\trigger{动手算之前，先问输出是否曾让输入变得可疑。}

\block{无噪二元信道}

\begin{example}[无噪二元信道]\label{ex:noiseless-binary}
设 $\X=\Y=\{0,1\}$，每个输入都被原样复现。
\begin{center}
\begin{tikzpicture}[x=2.4cm,y=1.1cm,>=latex,font=\small]
  \node (x0) at (0,1) {$0$};
  \node (x1) at (0,0) {$1$};
  \node (y0) at (1,1) {$0$};
  \node (y1) at (1,0) {$1$};
  \draw[->,indigo] (x0) -- (y0);
  \draw[->,indigo] (x1) -- (y1);
\end{tikzpicture}
\end{center}
每次使用信道恰好传送一个无差错比特，所以容量应当为 $1$。定义也给出同样的结论：
此时 $Y=X$，故 $\Ent(Y\mid X)=0$，于是
\[
  C=\max_{p(x)}\MI(X;Y)
   =\max_{p(x)}\Ent(X)
   =\answerin[1.2cm]{1}\ \text{比特},
\]
最大值在 $p(x)=\answerin[2.4cm]{(\tfrac12,\tfrac12)}$ 时取得。
\studentrecall{$C=1$ 比特，由均匀输入 $(\tfrac12,\tfrac12)$ 达到。}
\end{example}

\block{输出互不重叠的噪声信道}

\begin{example}[输出互不重叠]\label{ex:nonoverlapping}
设 $\X=\{0,1\}$，$\Y=\{1,2,3,4\}$，且
\[
  p(1\mid0)=p(2\mid0)=\tfrac12,
  \qquad
  p(3\mid1)=\tfrac13,\quad p(4\mid1)=\tfrac23 .
\]
\begin{center}
\begin{tikzpicture}[x=2.6cm,y=0.95cm,>=latex,font=\small]
  \node (x0) at (0,2.6) {$0$};
  \node (x1) at (0,0.4) {$1$};
  \node (y1) at (1,3.3) {$1$};
  \node (y2) at (1,2.1) {$2$};
  \node (y3) at (1,0.9) {$3$};
  \node (y4) at (1,-0.3) {$4$};
  \draw[->,indigo] (x0) -- node[above,sloped,font=\scriptsize]{$1/2$} (y1);
  \draw[->,indigo] (x0) -- node[below,sloped,font=\scriptsize]{$1/2$} (y2);
  \draw[->,madder] (x1) -- node[above,sloped,font=\scriptsize]{$1/3$} (y3);
  \draw[->,madder] (x1) -- node[below,sloped,font=\scriptsize]{$2/3$} (y4);
\end{tikzpicture}
\end{center}
这个信道确实有噪：两个输入的输出都是随机的。但输出集合 $\{1,2\}$ 与 $\{3,4\}$
\keyword{互不相交}，所以接收端从 $y$ 总能判断出送入的是哪个输入。于是
\[
  \Ent(X\mid Y)=\answerin[1.2cm]{0},
  \qquad
  \MI(X;Y)=\Ent(X)-\Ent(X\mid Y)=\Ent(X),
\]
\studentrecall{输出集合不相交意味着 $Y$ 决定 $X$，故 $\Ent(X\mid Y)=0$、$\MI=\Ent(X)$。}
从而
\[
  C=\max_{p(x)}\Ent(X)=1\ \text{比特},
\]
同样由均匀输入达到。每行内部的随机性无关紧要；只有行与行之间的重叠才消耗容量。
\end{example}

\begin{remark}
值得把一般原理点明：信道容量只通过“哪些输入可能被\emph{混淆}”依赖于转移矩阵。
若不同行的支撑两两不交，则该信道与 $|\X|$ 个符号上的无噪信道一样好，$C=\log|\X|$。
\end{remark}

\begin{yourturn}
\keyword{噪声打字机。}设 $\X=\Y=\{A,B,\ldots,Z\}$（26 个字母），每个字母以概率
$1/2$ 被正确打印、以概率 $1/2$ 被打印成下一个字母。若只用字母表中每隔一个的字母作
输入，则被使用输入的输出集合互不相交。这提示容量是多少？为什么不能更高？
\studentrecall{有 $13$ 个互不混淆的可用输入，$C=\log13=\log26-1$ 比特。}
\ifshowanswers\else\workspace[3]\fi
\answerprose{取相隔两位的 $13$ 个字母可使输出集合互不相交，这 $13$ 个输入构成一个
无噪信道，故 $C\ge\log13$。反之每行熵都是 $\Ent(Y\mid X=x)=1$ 比特，所以
$\MI(X;Y)=\Ent(Y)-1\le\log26-1=\log13$。因此 $C=\log13$ 比特；在全部 $26$ 个输入
上取均匀分布同样能达到。}
\end{yourturn}

\recall{真正毁掉容量的是行内的随机性，还是行与行之间的重叠？}


\clearpage
% ===========================================================================
\section{二元对称信道}\label{sec:bsc}
\warmth{0}\quad\whisper{最简单的含错信道，却已经包含了几乎全部困难。}

\begin{strand}
二元对称信道把每个输入比特以概率 $p$ 取反。它是含错信道最简单的模型，却抓住了
一般问题的大部分复杂性。其容量为 $1-h(p)$：每次使用一个比特，减去噪声注入的熵。
\end{strand}

\trigger{只要转移矩阵每一行的熵都相同，就用 $\MI=\Ent(Y)-\Ent(Y\mid X)$。}

\block{信道}

\begin{definition}[二元对称信道]\label{def:bsc}
\keyword{二元对称信道} $\mathrm{BSC}(p)$ 满足 $\X=\Y=\{0,1\}$，转移矩阵为
\[
  p(y\mid x)=
  \begin{pmatrix}
    1-p & p\\
    p & 1-p
  \end{pmatrix},
\]
即输入符号以概率 $p$ 被取反。
\begin{center}
\begin{tikzpicture}[x=2.8cm,y=1.5cm,>=latex,font=\small]
  \node (x0) at (0,1) {$0$};
  \node (x1) at (0,0) {$1$};
  \node (y0) at (1,1) {$0$};
  \node (y1) at (1,0) {$1$};
  \draw[->,indigo] (x0) -- node[above,font=\scriptsize]{$1-p$} (y0);
  \draw[->,indigo] (x1) -- node[below,font=\scriptsize]{$1-p$} (y1);
  \draw[->,madder] (x0) -- node[pos=0.22,above,sloped,font=\scriptsize]{$p$} (y1);
  \draw[->,madder] (x1) -- node[pos=0.22,below,sloped,font=\scriptsize]{$p$} (y0);
\end{tikzpicture}
\end{center}
两行都是向量 $(1-p,p)$ 的重排，故每一行的熵都等于
$\answerin[1.4cm]{h(p)}$。
\studentrecall{每一行都是 $(1-p,p)$ 分布，其熵为二元熵 $h(p)$。}
\end{definition}

\block{容量}

\begin{theorem}[BSC 的容量]\label{thm:bsc-capacity}
$\mathrm{BSC}(p)$ 的容量为
\[
  C=1-h(p)\quad\text{比特每次信道使用},
\]
由均匀输入分布 $p(x)=(\tfrac12,\tfrac12)$ 达到。
\end{theorem}

\begin{proof}
对任意输入分布，
\[
\begin{aligned}
  \MI(X;Y)
  &=\Ent(Y)-\Ent(Y\mid X)\\
  &=\Ent(Y)-\sum_{x}p(x)\,\Ent(Y\mid X=x)\\
  &=\Ent(Y)-\sum_{x}p(x)\,h(p)\\
  &=\Ent(Y)-h(p)\\
  &\le 1-h(p),
\end{aligned}
\]
最后一步是因为 $Y$ 是二元随机变量，故 $\Ent(Y)\le\log|\Y|=1$。取等要求 $Y$ 均匀。
取 $p(x)=(\tfrac12,\tfrac12)$ 时
\[
  \PP(Y=0)=\tfrac12(1-p)+\tfrac12p=\tfrac12,
\]
于是 $\Ent(Y)=1$，界被达到。因此 $C=1-h(p)$。
\end{proof}

\begin{remark}
让计算变短的关键一步是 $\Ent(Y\mid X)=h(p)$ \emph{与 $p(x)$ 无关}：权重 $p(x)$
乘的是彼此相同的行熵。这正是第 \ref{sec:symmetric} 节要推广的对称性。
\end{remark}
\loose{此处猜到均匀输入最优全靠对称性；没有对称性时见第 \ref{sec:properties} 节。}

\block{读懂这个答案}

\noindent\begin{tabularx}{\linewidth}{@{}l l L@{}}
\toprule
{\color{indigo}$p$} & {\color{indigo}$C=1-h(p)$} & \multicolumn{1}{l}{\color{madder}含义}\\
\midrule
$0$ & $1$ & 无噪信道；每次使用一个干净比特\\
$1$ & $1$ & 每个比特都被翻转，接收端直接翻回来即可\\
$1/2$ & $0$ & 输出与输入独立；什么也传不过去\\
$0.11$ & $\approx0.5$ & 大约每两次使用换来一个可用比特\\
\bottomrule
\end{tabularx}

\begin{yourturn}
用一句话解释为什么 $\mathrm{BSC}(1)$ 与 $\mathrm{BSC}(0)$ 容量相同，以及为什么
$p=\tfrac12$ 是最坏情形。
\studentrecall{确定性的翻转是可逆的；$p=\tfrac12$ 让 $Y$ 与 $X$ 独立。}
\ifshowanswers\else\workspace[2]\fi
\answerprose{$p=1$ 时翻转是确定性的，把输出重新贴标签即得无噪信道，且 $h(1)=0$。
$p=\tfrac12$ 时无论送什么，两个输出等可能，故 $Y$ 与 $X$ 独立，$h(\tfrac12)=1$、
$C=0$。}
\end{yourturn}

\begin{yourturn}
复述证明中的五行链条，并说明每一步的理由。
\studentrecall{拆开 $\MI$；把行熵取平均；各行熵都是 $h(p)$；用 $\Ent(Y)\le1$；均匀输入让 $Y$ 均匀。}
\ifshowanswers\else\workspace[4]\fi
\answerprose{$\MI=\Ent(Y)-\Ent(Y\mid X)$（定义）；展开
$\Ent(Y\mid X)=\sum_xp(x)\Ent(Y\mid X=x)$（条件熵定义）；每个
$\Ent(Y\mid X=x)=h(p)$（来自矩阵各行）；由 $\sum_xp(x)=1$ 该和为 $h(p)$；最后因
$Y$ 二元有 $\Ent(Y)\le1$，取等当且仅当 $Y$ 均匀，而均匀输入正好做到这一点。}
\end{yourturn}

\recall{为什么 BSC 的 $\Ent(Y\mid X)$ 不随 $p(x)$ 变化？}


\clearpage
% ===========================================================================
\section{二元删除信道}\label{sec:bec}
\warmth{0}\quad\whisper{丢掉一个比特，比被这个比特欺骗要便宜。}

\begin{strand}
删除信道会丢掉比例为 $\alpha$ 的比特，但从不篡改任何比特：接收端总是知道\emph{哪些}
比特丢了。它的容量是 $1-\alpha$——恰好是幸存下来的那部分信道使用。与 BSC 对照：
在 BSC 里翻转概率 $\alpha$ 要付的代价是 $h(\alpha)$，要贵得多。
\end{strand}

\trigger{当某个输出符号意味着“毫无信息”时，就对“它是否出现”这一事件做条件分解。}

\block{信道}

\begin{definition}[二元删除信道]\label{def:bec}
\keyword{二元删除信道} $\mathrm{BEC}(\alpha)$ 满足 $\X=\{0,1\}$、
$\Y=\{0,1,e\}$，其中 $e$ 是\keyword{删除}符号，且
\[
  p(0\mid0)=p(1\mid1)=1-\alpha,
  \qquad
  p(e\mid0)=p(e\mid1)=\alpha .
\]
\begin{center}
\begin{tikzpicture}[x=2.8cm,y=1.2cm,>=latex,font=\small]
  \node (x0) at (0,2) {$0$};
  \node (x1) at (0,0) {$1$};
  \node (y0) at (1,2) {$0$};
  \node (ye) at (1,1) {$e$};
  \node (y1) at (1,0) {$1$};
  \draw[->,indigo] (x0) -- node[above,font=\scriptsize]{$1-\alpha$} (y0);
  \draw[->,indigo] (x1) -- node[below,font=\scriptsize]{$1-\alpha$} (y1);
  \draw[->,madder] (x0) -- node[pos=0.6,above,sloped,font=\scriptsize]{$\alpha$} (ye);
  \draw[->,madder] (x1) -- node[pos=0.6,below,sloped,font=\scriptsize]{$\alpha$} (ye);
\end{tikzpicture}
\end{center}
任何比特都不会被取反，所以落在 $\{0,1\}$ 中的输出可确定地识别输入；输出为 $e$ 则
什么也没说。每行都是 $(1-\alpha,\alpha,0)$ 的重排，故
$\Ent(Y\mid X)=\answerin[1.6cm]{h(\alpha)}$。
\studentrecall{每行实际上是两点上的 $(1-\alpha,\alpha)$，行熵为 $h(\alpha)$。}
\end{definition}

\block{容量}

\begin{theorem}[BEC 的容量]\label{thm:bec-capacity}
$\mathrm{BEC}(\alpha)$ 的容量为
\[
  C=1-\alpha\quad\text{比特每次信道使用},
\]
由均匀输入分布达到。
\end{theorem}

\begin{proof}
照例写
\[
  C=\max_{p(x)}\MI(X;Y)
   =\max_{p(x)}\bigl[\Ent(Y)-\Ent(Y\mid X)\bigr]
   =\max_{p(x)}\Ent(Y)-h(\alpha).
\]
为计算 $\Ent(Y)$，令 $E=\mathbf{1}\{Y=e\}$ 为“发生删除”的指示变量。由于 $E$ 是
$Y$ 的函数，有 $\Ent(Y,E)=\Ent(Y)$，再用链式法则得
\[
  \Ent(Y)=\Ent(Y,E)=\Ent(E)+\Ent(Y\mid E).
\]
现在 $E\sim\operatorname{Bernoulli}(\alpha)$，故 $\Ent(E)=h(\alpha)$。给定
$E=1$ 时输出必为 $e$，贡献为零；给定 $E=0$ 时输出等于输入。记
$\PP(X=1)=\pi$，则
\[
  \Ent(Y\mid E)=\alpha\cdot0+(1-\alpha)\,h(\pi),
  \qquad\text{于是}\qquad
  \Ent(Y)=h(\alpha)+(1-\alpha)h(\pi).
\]
因此
\[
  C=\max_{\pi}\bigl[h(\alpha)+(1-\alpha)h(\pi)\bigr]-h(\alpha)
   =\max_{\pi}(1-\alpha)h(\pi)
   =1-\alpha,
\]
容量在 $\pi=\tfrac12$ 处达到。
\end{proof}

\begin{remark}
这里的解释是精确的而非近似的：发送的比特中有比例 $\alpha$ 被丢失，所以最多只有
比例 $1-\alpha$ 能通过，而这么多确实可以达到。还要注意两个“显然”的界在此表现不同：
$\Ent(Y)\le\log3$ \emph{不}是紧的，因为删除符号出现的频率永远不会超过 $\alpha$。
用删除指示变量拆分 $\Ent(Y)$，正是用来替代那个粗糙界的做法。
\end{remark}

\block{同一噪声水平下的 BSC 与 BEC}

\noindent\begin{tabularx}{\linewidth}{@{}l l L@{}}
\toprule
{\color{indigo}信道} & {\color{indigo}容量} & \multicolumn{1}{l}{\color{madder}接收端知道什么}\\
\midrule
$\mathrm{BSC}(p)$ & $1-h(p)$ & 收到的每个比特都可能是假的；错误位置未知\\
$\mathrm{BEC}(\alpha)$ & $1-\alpha$ & 收到的比特一定正确；删除位置被标记出来\\
\bottomrule
\end{tabularx}

取 $p=\alpha=0.1$ 时二者分别为 $C\approx0.531$ 与 $C=0.9$：知道损坏发生在\emph{哪里}
非常值钱。

\begin{yourturn}
关键技巧是 $\Ent(Y)=\Ent(E)+\Ent(Y\mid E)$。说出使这一等式在此既合法又有用的两个
事实。
\studentrecall{$E$ 是 $Y$ 的函数，故 $\Ent(Y,E)=\Ent(Y)$；给定 $E$ 后输出要么固定为 $e$，要么等于 $X$。}
\ifshowanswers\else\workspace[3]\fi
\answerprose{第一，$E=\mathbf{1}\{Y=e\}$ 是 $Y$ 的确定性函数，故
$\Ent(Y,E)=\Ent(Y)$，可以对这一对用链式法则。第二，对 $E$ 取条件把输出分成退化
情形（$Y=e$，熵为 $0$）与无噪情形（$Y=X$，熵为 $h(\pi)$），因此
$\Ent(Y\mid E)=(1-\alpha)h(\pi)$。}
\end{yourturn}

\begin{yourturn}
凭记忆补全两个容量：$\mathrm{BSC}(p)$ 的 $C=\answerin[2.2cm]{1-h(p)}$，
$\mathrm{BEC}(\alpha)$ 的 $C=\answerin[2.0cm]{1-\alpha}$。
\studentrecall{BSC 为 $1-h(p)$；BEC 为 $1-\alpha$。}
\end{yourturn}

\recall{为什么 $\Ent(Y)\le\log3$ 对 BEC 没有用？}


\clearpage
% ===========================================================================
\section{对称信道}\label{sec:symmetric}
\warmth{0}\quad\whisper{各行长得一样时，均匀输入最优，容量只是一行公式。}

\begin{strand}
BSC 的计算只用了两个事实：转移矩阵所有行的熵相同，以及均匀输入产生均匀输出。
具备这两点的信道称为（弱）对称信道；对它们而言容量就是 $\log|\Y|-\Ent(r)$，其中
$r$ 是任意一行——不需要做最大化。
\end{strand}

\trigger{先检查对称性；若信道对称，容量只花一行而不是一次最大化。}

\block{定义}

\begin{definition}[对称信道]\label{def:symmetric-channel}
若离散无记忆信道的转移矩阵 $p(y\mid x)$ 各行互为置换，且各\answerin[2.0cm]{列}
也互为置换，则称该信道\keyword{对称}。
\studentrecall{行互为置换，\emph{并且}列也互为置换。}
\end{definition}

\begin{definition}[弱对称信道]\label{def:weakly-symmetric-channel}
若 $p(y\mid x)$ 各行互为置换，且每一列的\answerin[1.6cm]{和}
$\sum_{x}p(y\mid x)$ 都相同，则称该信道\keyword{弱对称}。
\studentrecall{行互为置换，且所有列和相等。}
由于置换后的列有相同的和，每个对称信道都是弱对称的。
\end{definition}

\begin{example}[一个对称信道]\label{ex:symmetric-3x3}
设 $\X=\Y=\{1,2,3\}$，
\[
  p(y\mid x)=
  \begin{pmatrix}
    0.3 & 0.2 & 0.5\\
    0.5 & 0.3 & 0.2\\
    0.2 & 0.5 & 0.3
  \end{pmatrix}.
\]
每行都是 $r=(0.3,0.2,0.5)$ 的置换，每列也是同一向量的置换，故该信道对称。
\end{example}

\block{弱对称信道的容量}

\begin{theorem}[由对称性给出容量]\label{thm:symmetric-capacity}
对行向量为 $r$（转移矩阵的任意一行）的弱对称信道，
\[
  C=\log|\Y|-\Ent(r),
\]
且容量由 $\X$ 上的均匀分布达到。
\end{theorem}

\begin{proof}
由于各行都是 $r$ 的置换，每个行熵都等于 $\Ent(r)$，故对任意输入分布
\[
  \MI(X;Y)=\Ent(Y)-\Ent(Y\mid X)
          =\Ent(Y)-\sum_xp(x)\Ent(r)
          =\Ent(Y)-\Ent(r)
          \le\log|\Y|-\Ent(r),
\]
取等当且仅当 $Y$ 在 $\Y$ 上均匀。剩下要证均匀输入使 $Y$ 均匀。令
$c=\sum_xp(y\mid x)$ 为公共列和。取 $p(x)=1/|\X|$，则
\[
  p(y)=\sum_{x}\frac{1}{|\X|}\,p(y\mid x)=\frac{c}{|\X|},
\]
它不依赖 $y$；既然是概率分布，就必须等于 $1/|\Y|$。于是界被达到，
$C=\log|\Y|-\Ent(r)$。
\end{proof}

\begin{example}[例 \ref{ex:symmetric-3x3} 的容量]\label{ex:symmetric-3x3-capacity}
取 $r=(0.3,0.2,0.5)$，
\[
  \Ent(r)=-0.3\log0.3-0.2\log0.2-0.5\log0.5\approx1.485,
  \qquad
  \log|\Y|=\log3\approx1.585,
\]
故
\[
  C=\log 3-\Ent(r)\approx\answerin[1.8cm]{0.0995}\ \text{比特},
\]
由均匀输入 $p(x)=(\tfrac13,\tfrac13,\tfrac13)$ 达到。
\studentrecall{$C=\log3-\Ent(0.3,0.2,0.5)\approx1.585-1.485\approx0.10$ 比特。}
这个信道几乎没用：它的各行接近均匀，所以输出几乎不依赖输入。
\end{example}

\block{谁对称，谁不对称}

\noindent\begin{tabularx}{\linewidth}{@{}l l L@{}}
\toprule
{\color{indigo}信道} & {\color{indigo}对称？} & \multicolumn{1}{l}{\color{madder}用公式算容量}\\
\midrule
无噪二元信道 & 是 & $\log2-\Ent(1,0)=1$\\
$\mathrm{BSC}(p)$ & 是 & $\log2-h(p)=1-h(p)$\\
例 \ref{ex:symmetric-3x3} & 是 & $\log3-\Ent(r)\approx0.0995$\\
$\mathrm{BEC}(\alpha)$ & 否 & 列和为 $1-\alpha,1-\alpha,2\alpha$ 不等；公式不适用\\
\bottomrule
\end{tabularx}

BEC 的各行\emph{确实}互为置换，可见单靠行条件还不够：列和不相等时，均匀输入未必
产生均匀输出；事实上真实答案 $1-\alpha$ 与 $\log3-h(\alpha)$ 并不相同。

\begin{yourturn}
指出证明中用到对称性的两处，各对应一个定义条件。
\studentrecall{行互为置换使 $\Ent(Y\mid X)=\Ent(r)$ 为常数；列和相等使均匀输入产生均匀输出。}
\ifshowanswers\else\workspace[3]\fi
\answerprose{行互为置换使每个 $\Ent(Y\mid X=x)$ 都等于 $\Ent(r)$，于是无论 $p(x)$
如何都有 $\Ent(Y\mid X)=\Ent(r)$。列和相等使均匀输入下 $p(y)$ 为常数，从而
$\Ent(Y)=\log|\Y|$ 可以达到，不等式变成等式。}
\end{yourturn}

\begin{yourturn}
设 $\X=\Y=\{1,2,3\}$，每行都是 $(\tfrac12,\tfrac14,\tfrac14)$ 的置换且列也互为
置换，求该信道的容量。
\studentrecall{$C=\log3-\Ent(\tfrac12,\tfrac14,\tfrac14)=\log3-1.5\approx0.085$ 比特。}
\ifshowanswers\else\workspace[2]\fi
\answerprose{$\Ent(\tfrac12,\tfrac14,\tfrac14)=\tfrac12\cdot1+\tfrac14\cdot2+\tfrac14\cdot2=1.5$
比特，故 $C=\log3-1.5\approx1.585-1.5=0.085$ 比特。}
\end{yourturn}

\recall{BEC 违反了哪个条件？硬套公式会出什么问题？}


\clearpage
% ===========================================================================
\section{容量的性质}\label{sec:properties}
\warmth{0}\quad\whisper{单纯形上的凹函数最大化：最大值存在，而且有检验条件。}

\begin{strand}
容易的例子做完后还剩两个问题：$C$ 可以取哪些值？没有对称性直接送来最优输入时怎么办？
答案是：$C$ 介于 $0$ 与 $\log\min(|\X|,|\Y|)$ 之间；并且 $\MI(X;Y)$ 是 $p(x)$ 的
凹函数，所以这是一个性质良好的凸优化问题，且有可验证的最优性条件。
\end{strand}

\trigger{已经猜出候选最优输入、想证明它最优时，用等散度检验。}

\block{取值范围}

\begin{proposition}[初等界]\label{prop:capacity-bounds}
对任意离散无记忆信道，
\[
  0\le C\le\log\min\bigl(|\X|,|\Y|\bigr).
\]
并且 $C=0$ 当且仅当对任何输入分布 $X$ 与 $Y$ 都独立，即 $p(y\mid x)$ 各行全部相同。
\end{proposition}

\begin{proof}
非负性成立是因为总有 $\MI(X;Y)\ge0$，故最大值至少为 $0$。至于上界，
\[
  \MI(X;Y)\le\Ent(X)\le\log|\X|,
  \qquad
  \MI(X;Y)\le\Ent(Y)\le\log|\Y|,
\]
这里用了 $\MI(X;Y)=\Ent(X)-\Ent(X\mid Y)\le\Ent(X)$ 与对称的恒等式。取两个界中较小
者即得结论。若各行全部相同，则对每个 $x$ 都有 $p(y)=p(y\mid x)$，于是 $Y$ 与 $X$
独立、所有互信息为零；反之 $C=0$ 迫使支撑为整个 $\X$ 的均匀输入满足
$\MI(X;Y)=0$，从而对一切 $x$ 有 $p(y\mid x)=p(y)$。
\end{proof}

\begin{yourturn}
本讲的哪个信道分别达到两个极端 $C=0$ 与 $C=\log\min(|\X|,|\Y|)$？
\studentrecall{$\mathrm{BSC}(1/2)$ 给出 $C=0$；无噪二元信道与输出不重叠的信道给出 $C=1=\log2$。}
\ifshowanswers\else\workspace[2]\fi
\answerprose{$\mathrm{BSC}(\tfrac12)$ 的两行都是 $(\tfrac12,\tfrac12)$，故 $C=0$。
上界极端 $C=\log2=1$ 由无噪二元信道达到，也由输出互不重叠的信道达到，后者
$|\X|=2$ 且 $\Ent(X\mid Y)=0$。}
\end{yourturn}

\block{优化问题的形状}

\begin{proposition}[关于输入的凹性]\label{prop:concavity}
固定转移分布 $p(y\mid x)$，则映射 $p(x)\mapsto\MI(X;Y)$ 在 $\X$ 上的概率单纯形上
连续且\keyword{凹}。
\end{proposition}

\begin{proof}
写 $\MI(X;Y)=\Ent(Y)-\Ent(Y\mid X)$。第二项是
$\sum_xp(x)\Ent(Y\mid X=x)$，关于 $p(x)$ 线性。第一项是 $\Ent$ 在
$p(y)=\sum_xp(x)p(y\mid x)$ 处的取值，而这是 $p(x)$ 的线性映射；由于 $\Ent$ 是凹
函数，凹函数与线性映射的复合仍凹，故 $p(x)\mapsto\Ent(Y)$ 是凹的。凹函数减去线性
函数仍是凹的。连续性由公式即明。
\end{proof}

由此得到的都是实用推论：单纯形是紧集，故最大值可达；每个局部最大都是全局最大；
标准凸优化工具（或 Blahut--Arimoto 算法）可以数值计算 $C$。

\begin{remark}
与信源编码作对比：那里相关的凸性是关于被编码的分布的。这里分布反而是\emph{设计
变量}，信道是固定的。互信息在 $p(y\mid x)$ 固定时是 $p(x)$ 的凹函数，在 $p(x)$
固定时是 $p(y\mid x)$ 的凸函数。
\end{remark}

\block{一个最优性检验}

\begin{theorem}[等散度条件]\label{thm:kt-condition}
输入分布 $p^*(x)$ 及其诱导的输出分布 $p^*(y)$ 达到容量，当且仅当存在常数 $C$ 使得
\[
  \KL\bigl(p(y\mid x)\,\big\|\,p^*(y)\bigr)
  \begin{cases}
    =C, & \text{对每个 }p^*(x)>0\text{ 的 }x,\\[2pt]
    \le C, & \text{对每个 }p^*(x)=0\text{ 的 }x,
  \end{cases}
\]
此时该常数就等于容量。
\end{theorem}

背后的恒等式是
$\MI(X;Y)=\sum_xp(x)\,\KL\bigl(p(y\mid x)\|p(y)\bigr)$：互信息是信道各行与输出分布
之间散度的加权平均。若两个被使用的输入散度不同，把权重移向较大的一个就能提高这个
平均值——所以在最优点处所有被使用的行必须\answerin[3.6cm]{与 $p^*(y)$ 等距}。
\studentrecall{最优时每个被使用的输入行与输出分布的散度都相同，且等于 $C$。}

\begin{example}[核对 BEC 的答案]\label{ex:bec-check}
对 $\mathrm{BEC}(\alpha)$，均匀输入给出 $\{0,1,e\}$ 上的
$p^*(y)=\bigl(\tfrac{1-\alpha}{2},\tfrac{1-\alpha}{2},\alpha\bigr)$。对输入
$x=0$（其行为 $(1-\alpha,0,\alpha)$），
\[
  \KL\bigl(p(y\mid0)\,\big\|\,p^*(y)\bigr)
  =(1-\alpha)\log\frac{1-\alpha}{(1-\alpha)/2}
   +\alpha\log\frac{\alpha}{\alpha}
  =1-\alpha,
\]
由对称性，$x=1$ 同样如此。两个被使用的输入散度都是 $1-\alpha$，这同时验证了均匀
输入的最优性与容量值 $C=1-\alpha$。
\end{example}

\begin{yourturn}
对 $\mathrm{BSC}(p)$ 做同样的核对：均匀输入下输出均匀，计算
$\KL\bigl((1-p,p)\,\big\|\,(\tfrac12,\tfrac12)\bigr)$ 并与 $1-h(p)$ 比较。
\studentrecall{$\KL((1-p,p)\|(\tfrac12,\tfrac12))=1-h(p)$，两行结果相同。}
\ifshowanswers\else\workspace[3]\fi
\answerprose{$\KL\bigl((1-p,p)\|(\tfrac12,\tfrac12)\bigr)
=(1-p)\log\frac{1-p}{1/2}+p\log\frac{p}{1/2}
=1+(1-p)\log(1-p)+p\log p=1-h(p)$，$x=1$ 那一行完全相同。故均匀输入最优且
$C=1-h(p)$，与直接计算一致。}
\end{yourturn}

\loose{非对称信道（如只有一个输入会被干扰的 $Z$ 信道）最优输入不均匀——正好拿来练这个检验。}

\block{总结卡片}

\noindent\begin{tabularx}{\linewidth}{@{}l l L@{}}
\toprule
{\color{indigo}信道} & {\color{indigo}$C$} & \multicolumn{1}{l}{\color{madder}最优输入}\\
\midrule
无噪二元信道 & $1$ & 均匀\\
输出互不重叠 & $\log|\X|$ & 均匀\\
噪声打字机（$26$ 字母） & $\log13$ & 均匀\\
$\mathrm{BSC}(p)$ & $1-h(p)$ & 均匀\\
$\mathrm{BEC}(\alpha)$ & $1-\alpha$ & 均匀\\
弱对称，行为 $r$ & $\log|\Y|-\Ent(r)$ & 均匀\\
一般情形 & $\max_{p(x)}\MI(X;Y)$ & 凹最大化；等散度检验\\
\bottomrule
\end{tabularx}

上表中每一行的最优输入都是均匀的，因为这些信道都带有某种对称性；一般问题之难正是
来自不对称。

\recall{哪个恒等式把 $\MI(X;Y)$ 变成散度的平均？它在最优点说了什么？}


\end{document}
